% Einheit 5 von 5 – Drei Variablen und Lösungsvielfalt (Sek II) (bank/lineare-gleichungssysteme/e5.jsonl, zusammenbau v0.1)
%% TODO Zeitmarke Einheit 5: Marken-Zeile nicht lesbar
%% TODO Zweigzeile Teil 1: was hier gelernt wird (Satz in Schülersprache aus dem Einheitstitel) – Einheit 5
\einheitenkopf[e5]{Einheit 5 von 5 · Drei Variablen und Lösungsvielfalt (Sek II)}
\zweigzeile{keine Prüfungsaufgabe}
\setcounter{aufgabe}{28}

%% TODO Titel: Ich-kann-Satz fehlt in der Bank (Platzhalter: Kettenname) – Nr. 29
%% TODO Anweisung: ein Satz über der ersten Teilaufgabe fehlt in der Bank – Nr. 29
\begin{aufgabe}{Drei Variablen und Lösungsvielfalt}
%% TODO \rechenplatz{n}: Schrittzahl des Grundfalls fehlt in der Bank (nur bei fester Schreibform, 2.3 b)
\begin{teile}
\teil Wie viele Lösungen? Beim Lösen eines Systems mit drei Variablen steht am Ende $0 = 0$. Kreuze an.\\ \kreuz{genau eine}\\ \kreuz{keine}\\ \kreuz{unendlich viele}
\teil Lösung nachweisen: Zeige durch Einsetzen, dass $x = 2$, $y = 1$, $z = 1$ das System I: $x + 2y - z = 3$, II: $3x + y - 2z = 5$ löst.
\end{teile}
\begin{gleichungsraster}[2]
\gl{\text{Gestaffelt lösen: I: $x + y - z = 2$, II: $2y + z = 7$, III: $3z = 9$ – Lösung? Löse von unten nach oben.}} &
\gl{\text{Drei Variablen: I: $2x + y = 9$, II: $y + 3z = 9$, III: $y + z = 5$ – Lösung? Subtrahiere zwei Gleichungen und setze ein.}} \\
\gl{\text{Wie viele Lösungen? I: $x + 2y + z = 4$, II: $x - y + 2z = 1$, III: $2x + y + 3z = 5$. Entscheide rechnerisch.}} &
\gl{\text{Schar und Auswahl: I: $x - z = 1$, II: $y + z = 5$, III: $3y + 3z = 15$. Gib die Lösungsmenge mit $z = t$ an und wähle die Lösung, in der $x$, $y$ und $z$ positiv und ganzzahlig sind und $x$ am größten ist.}} \\
\end{gleichungsraster}
\begin{teile}
\teil Erweitertes System: I: $x + y = 5$, II: $x - y = 1$ haben genau eine gemeinsame Lösung. Dazu kommt III: $2x + y = a$ mit einer reellen Zahl $a$. Wie viele Lösungen hat das System aus I, II und III, abhängig von $a$?
\teil Fallunterscheidung: I: $x + y + z = 4$, II: $y - z = 2$, III: $(a^2 - 9) \cdot z = a + 3$ mit einer reellen Zahl $a$. Wie viele Lösungen, abhängig von $a$?
\teil Aussage nachweisen: Eine Mischung aus drei Sorten hat die Anteile $x$, $y$ und $z$; keiner ist negativ. Das zugehörige Gleichungssystem hat die Lösungsmenge $\{(3 - 4t;\ -2 + 3t;\ t) \mid t \in \mathbb{R}\}$. Zeige: der Anteil $z$ ist mindestens doppelt so groß wie der Anteil $x$. \hfill (Abitur 2024 LK)
\end{teile}
\end{aufgabe}

%% TODO Titel: Ich-kann-Satz fehlt in der Bank (Platzhalter: Kettenname) – Nr. 30
%% TODO Anweisung: ein Satz über der ersten Teilaufgabe fehlt in der Bank – Nr. 30
\begin{aufgabe}{Drei Variablen und Lösungsvielfalt – Fehler finden, Begründen}
\begin{teile}
\teil Aus einem gestaffelten System bleibt III: $(a - 2) \cdot z = 6$. Jan rechnet \rechnung{z &= \frac{6}{a - 2}} und schreibt: „Das System hat für jedes $a$ genau eine Lösung.“ Finde den Fehler.
\teil Ole löst I: $x + y + z = 2$, II: $y + 2z = 5$, III: $z = -1$: \rechnung{\text{II}: y + 2 \cdot (-1) &= 5 \\ y &= 5 - 2 = 3 \\ \text{I}: x + 3 - 1 &= 2 \\ x &= 0} Finde den Fehler.
\teil Aus der Schar $x = \frac{t}{3}$, $y = 5 - t$, $z = t$ soll die Lösung mit lauter positiven ganzen Zahlen und dem größten $z$ gewählt werden. Eva nimmt $t = 4$ und erhält $x = \frac{4}{3}$, $y = 1$, $z = 4$. Finde den Fehler.
\teil Warum kann eine Gleichung $0 \cdot z = c$ keine oder unendlich viele Lösungen haben? Begründe.
\teil Warum liefert eine frei wählbare Variable $z = t$ unendlich viele Lösungen, und wie wählen Bedingungen wie „alle ganzzahlig und positiv“ daraus eine aus? Begründe.
\end{teile}
\end{aufgabe}
